% il existe plusieurs classes de documents
% pour des documents plus longs, vous pouvez utiliser
% book ou report
% remplacer oneside par twoside pour les documents recto-verso

\documentclass[a4paper, oneside]{article}

% vous pouvez changer les paramètres, voici les options dispoinibles :
% - a4paper
% - fancysections
% - notitlepage ou titlepage
% - sectionmark (pour avoir le rappel de la section dans les en-têtes)
% - chaptermark (de même pour les chapitres)
% - markboth (pour avoir la section et le chapitre dans les en-têtes)
% - pagenumber
% - enmanquedinspiration (pour des marges plus grandes)
% en cas de doutes, pas de doutes, la documentation est sur :
% https://gitlab.binets.fr/typographix/polytechnique/-/blob/master/source/polytechnique.pdf

\usepackage[a4paper,  fancysections,  titlepage]{polytechnique}
\usepackage[french]{babel}
\usepackage[T1]{fontenc}
\usepackage{blindtext}
\usepackage[hidelinks]{hyperref}
\usepackage{amsmath, amssymb}

% nous avons défini deux commandes :
\newcommand{\code}[1]{%
    \mbox{\ttfamily%
        \detokenize{#1}%
    }%
}

\newcommand{\resultat}[1]{%
    \quad \rightsquigarrow \quad #1%
}

\author{Louis Vaneau}
\date{\today}
\title{Grand titre}
\subtitle{Petit titre}

% pour changer de logo, ajoutez l'image dans un fomat PDF
% ou image en la glissant à droite et remplacez typographix
% par le nom de l'image, si vous ne voulez pas de logo,
% supprimez la ligne
\logo{typographix}


\begin{document}
	\maketitle 
	
	\tableofcontents
	
	\section{Ma première section}
	
	Un petit paragraphe introductif. On change de paragraphe en sautant deux lignes.
	
	Ceci est un autre paragraphe.
	
	
	\subsection{Une sous section}
	
	\subsection*{Une section sans numéro}
	
	\subsection{Faire des énumérations}
	Voici une liste :
	\begin{itemize}
		\item un premier élément ;
		\item un second élément.
	\end{itemize}
	
	On peut faire des listes numérotées :
	\begin{enumerate}
		\item de 1 ;
		\item de 2.
	\end{enumerate}
	
	\begin{description}
		\item [Éléphant :] gros mammifère.
        \item [Moineau :] petit oiseau.
	\end{description}

    \section{Des majuscules}
	Pour accentuer les majuscules, vous pouvez écrire \og É \fg{} ou \og \'E \fg{}.	On peut aussi faire référence à des articles scientifiques.\cite{premier_article}
	
	On peut mettre des liens comme celle-ci \href{https://typographix.binets.fr}{pour le site de typographix}.

 \section{Les mathématiques}
 Pour écrire des formules simples dans un paragraphe, on peut utiliser un mode \emph{inline math} : \(x^2+2x+1=0\). Pour des formules plus compliquées qui seraient toutes petites s'il fallait les mettre sur une ligne de texte normale, on peut utiliser le mode formules :
 \[
    \code{\sum_{k=0}^n u_k = \int_0^{+\infty} f(t) \: \mathrm{d} t} \resultat{\sum_{k=0}^n u_i = \int_0^{+\infty}f(t)\:\mathrm{d}t}
 \]

 Pour des mathématiques perfectionnistes quelques règles sont à observer : 
 \begin{itemize}
    \item la plupart des signes que l'on écrit dans l'environnement mathématique sont interprétés comme des variables (et sont donc en italique) ce qui n'est pas toujours souhaitable ;
    \item les signes comme \(+\) sont interprétés tantôt comme des opérateurs unaires ou binaires (selon le contexte), ce qui influe sur les espacements avant ou après ;
    \item les parenthèses sont en général bien interprétées comme des symboles d'ouverture ou de fermeture, mais ce n'est généralement pas le cas pour d'autres symboles. Par exemple :
    \begin{align*}
        \code{ ]-a , b] \neq \mathopen] - a, b \mathclose] } & \resultat{ ]-a , b] \neq \mathopen] - a, b \mathclose] }\\
        \code{ |\psi> = a|+> + b|-> } & \resultat{ |\psi> = a|+> + b|-> } \\
        \code{ \mathopen|\psi \mathclose> = \dots } & \resultat{ \mathopen|\psi \mathclose> = \dots }
    \end{align*}
    Dans le premier exemple, on voit que le signe $-$ peut être interprété comme un opérateur binaire ($] - a$) ou unaire ($\mathopen] - a$). Le deuxième exemple montre comment une équation peut devenir illisible à cause de ces erreurs d'attention ;
    \item les opérateurs mathématiques, comme les fonctions, doivent être définis comme tels, de manière à être affichées en romain, et à placer l'objet sur lequel il agit correctement, par exemple :
    \begin{equation*}
        \code{sin-x \neq \sin-x} \resultat{ sin-x \neq \sin-x }
    \end{equation*}
    
    Là encore, $-$ peut être interprété comme un opérateur unaire ou binaire. Si l'opérateur n'est pas déjà défini, on peut utiliser \code{\DeclareMathOperator{\grad}{grad}}, que l'on mettra dans le préambule (c'est-à-dire avant \code{\begin{document}}) ;
    \item les fonctions et les constantes sont en style romain : \(i\) est une variable quelconque, \(\mathrm i\) est un nombre complexe de module 1 et orthogonal à 1. De même, les $\mathrm d$ des éléments différentiels sont droits, ainsi que les constantes physiques. Cela vaut aussi pour les indices qui sont des abréviations comme l'énergie initiale \(E_{\mathrm i}\) où le $\mathrm i$ signifie initial.
\end{itemize}

Évidemment, vous êtes libre de respecter ou pas chacune de ses règles, certaines aident à la lisibilité, d'autres relèvent plus du détail. Pensez cependant que vos documents seront bien plus beaux si vous les respectez, et certains sauront remarquer le soin que vous y avez apporté.

Il faut surtout retenir que \LaTeX{} est un \emph{programme}, qui doit interpréter ce que vous écrivez : le résultat sera donc bien meilleur si vous utilisez les indications d'interprétation que sont \code{\mathopen} ou \code{\sin}.

\begin{thebibliography}{}
	\bibitem{premier_article}{Karen Sparck Jones. 1972. A Statistical Interpretation of Term Specificity and its Application in Retrieval. \emph{ Journal of Documentation.}Vol. 28 No. 1, pp. 11-21
	}
\end{thebibliography}
	
\end{document}